%% Elementary categorical definitions and complete proof-annex statements.

\subsection{Categories, functors, and natural transformations}
\label{sec:p10-categorical-toolkit}

\begin{hbconvention}{def:p10-categorical-size}{Size conventions for categorical constructions}
The fixed universe \(\mathcal U\) of small sets is transitive,
has \(\mathbb N\in\mathcal U\), and is closed under pairing, powersets,
and unions of small-indexed families:
\[
 \begin{gathered}
 a\in\mathcal U\ \Longrightarrow\ a\subseteq\mathcal U,\qquad
 a,b\in\mathcal U\ \Longrightarrow\ \{a,b\}\in\mathcal U,\\
 a\in\mathcal U\ \Longrightarrow\ \mathcal P(a)\in\mathcal U,\\
 I\in\mathcal U,\quad A_i\in\mathcal U\ (i\in I)
       \ \Longrightarrow\ \bigcup_{i\in I}A_i\in\mathcal U .
 \end{gathered}
\]
These closures also keep subsets, function sets, quotients, and
small-indexed products within \(\mathcal U\).
The ambient collection of all sets is a permitted default:
in that reading small means set-sized, with the ordinary
set/class distinction, and the closures are ordinary set-theoretic
properties.  If a set-sized \(\mathcal U\) is used for relativization,
it is required to be a Grothendieck universe with the stated
closures.  The existence of such a set universe, or a larger
containing one, is an additional size hypothesis when invoked;
no universe-existence assertion is part of the categorical proofs.

A category is small when its object and arrow sets are small.
A diagram is small when its indexing category is small; its
values may be objects of a large category.
Families indexed by a large category are handled in the containing
ambient framework.  Statements about class-sized natural
transformations use their displayed component families and
parametrizations; they do not assert a set of all such families.
A small set of parameters need not make every entire component
family a small object.
\end{hbconvention}

\begin{hbdefinition}{def:p10-category}{Category, functor, and natural transformation}
A category \(\mathcal C\) consists of objects, arrows with specified
domains and codomains, an identity \(1_A:A\to A\) for each object,
and a composite \(g\circ f:A\to C\) whenever \(f:A\to B\) and
\(g:B\to C\).  For composable arrows,
\[
 h\circ(g\circ f)=(h\circ g)\circ f,\qquad
 1_B\circ f=f=f\circ1_A.
\]
Write \(\mathcal C(A,B)\) for the arrows \(A\to B\).
The category is locally small if each such collection is a small
set, and small if its objects and all its arrows form small sets.
An isomorphism \(f:A\to B\) has an arrow \(g:B\to A\) with
\(g f=1_A\) and \(f g=1_B\).

A functor \(F:\mathcal C\to\mathcal D\) assigns objects and arrows
with the indicated domains and codomains, and satisfies
\[
 F(1_A)=1_{F A},\qquad F(g\circ f)=F(g)\circ F(f).
\]
The identity functor is \(1_{\mathcal C}\).
The opposite category \(\mathcal C^{\mathrm{op}}\) reverses every
arrow and the order of composition.  A contravariant set-valued
functor is an ordinary functor
\(F:\mathcal C^{\mathrm{op}}\to\mathsf{Set}\): thus \(f:A\to B\)
induces \(F(f):F(B)\to F(A)\).

For functors \(F,G:\mathcal C\to\mathcal D\), a natural
transformation \(\alpha:F\Rightarrow G\) consists of arrows
\(\alpha_A:F A\to G A\) satisfying
\[
                  G(f)\circ\alpha_A=\alpha_B\circ F(f)
                  \quad(f:A\to B).
\]
It is a natural isomorphism when every component is an
isomorphism.  The naturality equation then also holds for the
componentwise inverses.  Transformations compose componentwise;
identity components give the identity transformation.
\index{category}\index{functor}\index{natural transformation}
\index{opposite category}\index{locally small category}
\end{hbdefinition}

The category \(\mathsf{Set}\) has small sets as objects and functions
as arrows.  A small partially ordered set gives a small category with one
arrow \(a\to b\) when \(a\leq b\), and none otherwise.  Its identity
and composition laws are reflexivity and transitivity; antisymmetry
makes isomorphic objects equal.

\begin{hbdefinition}{def:p10-category-equivalence}{Fullness, faithfulness, and equivalence}
A functor \(F:\mathcal C\to\mathcal D\) is full when each map
\(\mathcal C(A,B)\to\mathcal D(F A,F B)\) is surjective, and faithful
when each is injective.  It is essentially surjective when every
object of \(\mathcal D\) is isomorphic to some \(F A\).
An equivalence consists of functors \(F:\mathcal C\to\mathcal D\)
and \(G:\mathcal D\to\mathcal C\), together with natural
isomorphisms \(1_{\mathcal C}\Rightarrow GF\) and
\(FG\Rightarrow1_{\mathcal D}\).
The functor \(G\) is a quasi-inverse.  The particular equivalence
constructed below supplies these functors and isomorphisms
explicitly.
\index{equivalence!of categories}\index{full functor}
\index{faithful functor}\index{essentially surjective functor}
\end{hbdefinition}

These definitions use the elementary conventions of
\cite[Chapter~1, especially Definitions~1.1.1, 1.3.1, 1.4.1,
and 1.4.3]{Riehl2026ContextOnline}.  The classical references
\cite{MacLane1998,Awodey2010} give separate book presentations.
The Riehl locators in this section refer to the author-hosted
second-edition text retrieved on 6 September 2026.

\subsection{Limits and colimits}

\begin{hbdefinition}{def:p10-limit}{Diagrams, cones, limits, and colimits}
Let \(J\) be a small category and \(D:J\to\mathcal C\) a functor,
called a diagram of shape \(J\).
A cone with vertex \(X\) is a family \(p_j:X\to D j\) satisfying
\[
                         D(u)\circ p_j=p_k
                         \quad(u:j\to k\text{ in }J).
\]
A limit is a cone \((L,(\ell_j)_j)\) such that every cone
\((X,(p_j)_j)\) has exactly one mediating arrow \(a:X\to L\)
with \(\ell_j a=p_j\) for all \(j\).
Equivalently, it is a terminal object in the category of cones,
whose arrows preserve every cone leg.

A cocone with vertex \(X\) has arrows \(q_j:D j\to X\) satisfying
\(q_k D(u)=q_j\).  A colimit is a cocone \((K,(k_j)_j)\) such that
every cocone \(q_j:D j\to X\) has exactly one arrow \(b:K\to X\)
with \(b k_j=q_j\).  It is an initial object in the category of
cocones.  The empty diagram has no legs or commutativity conditions.
Its limits are terminal objects and its colimits are initial
objects.
\index{limit!categorical}\index{colimit}\index{cone!categorical}
\index{terminal object}\index{initial object}
\end{hbdefinition}

\begin{hbdefinition}{def:p10-products-equalizers}{Products, equalizers, and their duals}
For a small family \((A_i)_{i\in I}\), a product is an object
\(\prod_i A_i\) with projections \(\pi_i\) for which every family
\(f_i:X\to A_i\) has a unique arrow
\(\langle f_i\rangle_i:X\to\prod_i A_i\) satisfying
\(\pi_i\langle f_i\rangle_i=f_i\).
It is the limit of the discrete diagram on \(I\).
A coproduct \(\coprod_i A_i\) has injections \(\iota_i\) and the
dual unique-arrow property for families \(A_i\to X\).
The empty product is terminal; the empty coproduct is initial.

For parallel arrows \(f,g:A\to B\), an equalizer is an arrow
\(e:E\to A\) with \(f e=g e\), universal among arrows with this
property: every \(a:X\to A\) with \(f a=g a\) factors uniquely
as \(a=e h\).
A coequalizer is \(q:B\to Q\) with \(q f=q g\), and every
\(b:B\to X\) with \(b f=b g\) factors uniquely as \(b=h q\).

A pullback of \(f:A\to C\) and \(g:B\to C\) consists of
\(p:P\to A\), \(q:P\to B\), with \(f p=g q\), such that every
commuting pair \(a:X\to A,b:X\to B\) factors uniquely through
\((p,q)\).  A pushout is the dual construction for
\(A\leftarrow C\to B\).  These are limits and colimits of the
indicated three-object diagrams.

An arrow \(m:A\to B\) is a monomorphism when
\(m r=m s\) implies \(r=s\) for every pair \(r,s:X\to A\).
An epimorphism satisfies the dual right-cancellation property.
\index{product!categorical}\index{coproduct}\index{equalizer}
\index{coequalizer}\index{pullback}\index{pushout}
\index{monomorphism}\index{epimorphism}
\end{hbdefinition}

For chosen binary products, write
\[
 f\times g=\langle f\pi_1,g\pi_2\rangle:
           A\times B\longrightarrow A'\times B'
 \quad(f:A\to A',\ g:B\to B').
\]
The projections distinguish arrows into a product.  Thus equations
between such arrows can be established by composing with both
projections.

\begin{hbtheorem}{thm:p10-limits-unique}{proved=annex:p10.cat.1}{Uniqueness of universal constructions}
Two limiting cones for the same diagram have a unique isomorphism
that preserves every cone leg.  Two colimiting cocones have the
dual unique isomorphism.
Every equalizer arrow is monic, and every coequalizer arrow is epic.
For products, pairing commutes with precomposition and product
arrows preserve identities and composition:
\[
 \langle f,g\rangle a=\langle f a,g a\rangle,\qquad
 1_A\times1_B=1_{A\times B},\qquad
 (f'\times g')(f\times g)=(f'f)\times(g'g).
\]
The uniqueness assertion concerns isomorphisms compatible with
the specified legs; it imposes no uniqueness on arbitrary
isomorphisms of the underlying objects.
\end{hbtheorem}

\begin{hbtheorem}{thm:p10-limits-products-equalizers}{proved=annex:p10.cat.2}{Constructing limits and colimits}
A locally small category with all small products and equalizers
has all small limits.  A locally small category with all small
coproducts and coequalizers has all small colimits.
Terminal objects, binary products, and equalizers suffice for
finite limits; initial objects, binary coproducts, and
coequalizers suffice for finite colimits.

For a small diagram \(D:J\to\mathsf{Set}\), a limit is the set
\[
 \left\{(x_j)_j\in\prod_{j\in\operatorname{Ob}J}D j:
       D(u)(x_j)=x_k\text{ for every }u:j\to k\right\}.
\]
A colimit is the disjoint union of the \(D j\)'s modulo the least
equivalence relation identifying
\((j,x)\) with \((k,D(u)(x))\) for every \(u:j\to k\).
The empty-diagram cases give a singleton and the empty set.
\end{hbtheorem}

Cones and their universal properties are presented in
\cite[Section~3.1, Definitions~3.1.1--3.1.6 and
Proposition~3.1.7]{Riehl2026ContextOnline}.
The general product/equalizer and coproduct/coequalizer
constructions are Theorem~3.5.11, pp.~112--113, of that text.

\subsection{Representable functors and Yoneda}

For an object \(A\) of a locally small category, the representable
functor \(\mathcal C(-,A):\mathcal C^{\mathrm{op}}\to\mathsf{Set}\)
sends \(X\) to \(\mathcal C(X,A)\), and \(u:Y\to X\) to
precomposition \(f\mapsto f u\).
Write \(\operatorname{Nat}(H,F)\) for natural transformations
between functors with the same domain and codomain.

\begin{hbtheorem}{thm:p10-yoneda}{proved=annex:p10.cat.3}{Yoneda lemma}
For a locally small category \(\mathcal C\), a functor
\(F:\mathcal C^{\mathrm{op}}\to\mathsf{Set}\), and an object \(A\),
evaluation at \(1_A\) gives a bijection
\[
 \operatorname{Nat}(\mathcal C(-,A),F)\longrightarrow F(A),
 \qquad \alpha\longmapsto\alpha_A(1_A).
\]
The inverse sends \(x\in F(A)\) to the transformation
\(\alpha^x_X(f)=F(f)(x)\).
This bijection is natural in \(A\) and \(F\).
It gives a canonical parametrization of the displayed family of
transformations by the small set \(F(A)\).
When \(\mathcal C\) is large, this does not assert that each
transformation's entire component family belongs to the fixed
universe.

Every natural transformation
\(\mathcal C(-,A)\Rightarrow\mathcal C(-,B)\) has the form
\(f\mapsto h f\) for a unique \(h:A\to B\).
Thus the Yoneda assignment \(A\mapsto\mathcal C(-,A)\) is full
and faithful on these hom-sets, and isomorphic representable
functors have isomorphic representing objects.
\index{Yoneda lemma}\index{representable functor}
\end{hbtheorem}

This is the contravariant form of
\cite[Theorem~2.2.4, pp.~61--63; Exercise~2.2.i,
p.~66]{Riehl2026ContextOnline}, with its proof supplied here.

\subsection{Adjunctions}

\begin{hbdefinition}{def:p10-adjunction}{Adjunction by natural hom-set bijections}
Let \(F:\mathcal C\to\mathcal D\) and \(G:\mathcal D\to\mathcal C\)
be functors between locally small categories.
An adjunction \(F\dashv G\) is a family of bijections
\[
 \Phi_{A,B}:\mathcal D(F A,B)\longrightarrow\mathcal C(A,G B)
\]
natural in both variables.  Explicitly, for
\(u:A'\to A\), \(f:F A\to B\), and \(v:B\to B'\),
\[
 \Phi_{A',B'}(v f F(u))=G(v)\,\Phi_{A,B}(f)\,u.
\]
The functor \(F\) is the left adjoint and \(G\) the right adjoint.
The image \(\Phi(f)\) and its inverse image are called
transposes with respect to this adjunction.
\index{adjunction}\index{left adjoint}\index{right adjoint}
\end{hbdefinition}

\begin{hbtheorem}{thm:p10-adjunction-data}{proved=annex:p10.cat.4}{Units, counits, and triangle identities}
Adjunction bijections \(\Phi\) determine natural transformations
\[
 \eta:1_{\mathcal C}\Rightarrow GF,\qquad
 \varepsilon:FG\Rightarrow1_{\mathcal D}
\]
by
\[
 \eta_A=\Phi_{A,F A}(1_{F A}),\qquad
 \varepsilon_B=\Phi^{-1}_{G B,B}(1_{G B}).
\]
They satisfy the triangle identities
\[
 \varepsilon_{F A}\circ F(\eta_A)=1_{F A},\qquad
 G(\varepsilon_B)\circ\eta_{G B}=1_{G B}.
\]
Conversely, natural transformations \(\eta,\varepsilon\) of these
types satisfying these identities determine an adjunction by the
mutually inverse maps
\[
 f\longmapsto G(f)\circ\eta_A,\qquad
 g\longmapsto\varepsilon_B\circ F(g).
\]
The constructions recover the original bijections and the original
unit and counit.
\index{unit!of an adjunction}\index{counit}
\index{triangle identities}
\end{hbtheorem}

\begin{hbtheorem}{thm:p10-adjoint-preservation}{proved=annex:p10.cat.5}{Adjoints preserve universal constructions}
If \(F\dashv G\), then \(F\) sends every existing small colimiting
cocone in \(\mathcal C\) to a colimiting cocone in \(\mathcal D\),
and \(G\) sends every existing small limiting cone in
\(\mathcal D\) to a limiting cone in \(\mathcal C\).
Consequently left adjoints preserve initial objects, coproducts,
coequalizers, and pushouts whenever these exist; right adjoints
preserve terminal objects, products, equalizers, and pullbacks.
The conclusion identifies the image cone or cocone by its
universal property, independently of any different chosen
representative of the same limit or colimit.
\end{hbtheorem}

The hom-set and unit/counit presentations are
\cite[Definition~4.1.1, p.~132; Definition~4.2.5 and
Theorem~4.2.7, pp.~140--142]{Riehl2026ContextOnline}.
Preservation of existing limits and colimits is Theorem~4.5.2,
pp.~158--159, of that text.

\subsection{Cartesian closed categories}

\begin{hbdefinition}{def:p10-ccc}{Cartesian closed category}
A cartesian closed category is a locally small category with a
terminal object, binary products, and exponentials.
When interpreting syntax, specify choices of these constructions.
For objects \(A,B\), an exponential consists of an object \(B^A\)
and an evaluation arrow
\[
                    \operatorname{ev}_{A,B}:B^A\times A\to B
\]
such that every \(f:X\times A\to B\) has a unique transpose
\(\Lambda(f):X\to B^A\) satisfying
\[
 \operatorname{ev}_{A,B}\circ(\Lambda(f)\times1_A)=f.
\]
Equivalently, the two maps
\[
 \mathcal C(X\times A,B)\ \longleftrightarrow\
 \mathcal C(X,B^A),\qquad
 f\longmapsto\Lambda(f),\quad
 k\longmapsto\operatorname{ev}_{A,B}\circ(k\times1_A)
\]
are inverse.  The inverse laws are the categorical \(\beta\)
and \(\eta\) equations.
In particular,
\(\Lambda(\operatorname{ev}_{A,B}\circ(k\times1_A))=k\).
\index{cartesian closed category}\index{exponential!categorical}
\index{evaluation!of an exponential}
\index{Curry--Howard--Lambek correspondence}
\end{hbdefinition}

\begin{hbtheorem}{thm:p10-exponential-adjunction}{proved=annex:p10.cat.6}{Exponential adjunction and examples}
In a cartesian closed category, transposition is natural in its
parameter:
\[
 \Lambda(f\circ(s\times1_A))=\Lambda(f)\circ s
       \quad(f:X\times A\to B,\ s:Y\to X).
\]
For fixed \(A\), the assignments \(X\mapsto X\times A\) and
\(B\mapsto B^A\) extend to functors with
\[
                         (-)\times A\dashv(-)^A.
\]
The category \(\mathsf{Set}\) is cartesian closed, with a singleton
as terminal object, Cartesian products as products, and the set
of functions \(A\to B\) as \(B^A\).
A Heyting algebra, regarded as a category under its order, is
cartesian closed with terminal object \(1\), product \(a\wedge b\),
and exponential \(b^a=(a\Rightarrow b)\).
\end{hbtheorem}

Function-type semantics uses precisely these evaluation and
transposition operations
\cite[Section~3.3, pp.~22--24]{Pitts1995Draft}.
The next construction applies them to the registered arrow-only
typed lambda calculus; the following internal-language construction
gives an explicit equivalence for a named small category.
